% chapters/testing/integration_tests.tex
% Sección 6.3: Pruebas de integración (~1 página)
% ============================================================================
%
% GUÍA: Describir las pruebas de integración de endpoints.
%       Referencia a \ref{rnf:test:integration}.
%
%   1. Qué se prueba:
%      - Flujo completo petición-respuesta de cada endpoint principal.
%      - Autenticación y permisos (usuario autorizado/no autorizado).
%      - Validaciones (entradas correctas e incorrectas).
%      - Códigos de respuesta HTTP esperados.
%      - Paginación, filtrado, ordenación.
%
%   2. Cómo se ejecutan:
%      - APITestCase de DRF / pytest-django.
%      - Base de datos transaccional (rollback tras cada test).
%      - Factories para generar datos de prueba.
%
%   3. Ejemplo representativo:
%      \PythonCode[cod:integration_test]{Test de integración de ejemplo}
%                 {Prueba de integración del endpoint de inicio de sesión}
%                 {codes/test_login_endpoint.py}{1}{25}{1}
%      TODO: Copiar un test de integración representativo a codes/.
%
%   4. Resultados resumidos:
%      - Número de tests de integración.
%      - Endpoints cubiertos / total de endpoints.
%      - Tabla resumen por recurso (similar a la de unitarias).
%
%   NOTA: Resultados detallados van en el \cref{ap:detailed_results}.
% ============================================================================

% TODO: Redactar la sección de pruebas de integración.
